\subsection{What the non-Ferrari affine quotient sees in the cubic residual}

\paragraph{D23.F27.}
In the literal non-Ferrari quotient $\eta_3$ of supplied FP.31,
\[
 a\mapsto(0,U),\quad b_0\mapsto(0,L),\quad
 c\mapsto(e_1,U),\qquad
 U=\begin{pmatrix}1&1\\0&1\end{pmatrix},\quad
 L=\begin{pmatrix}1&0\\-1&1\end{pmatrix},
\]
retain $\delta=ab_0$, $K=c^2$ and the cubic residual $D_h$ of
D23.C17.  For an actual gentle parameter, the exceptional value
$g=f(\delta,K)$ has the necessary form
\[
 \eta_3(g)=(u e_1,N_\rho),\quad u\in\mathbb F_3,\qquad
 N_0=I,\quad N_1=\begin{pmatrix}0&1\\2&0\end{pmatrix},\quad
 N_2=\begin{pmatrix}1&1\\1&2\end{pmatrix}.
\tag{D23.F27.1}
\]
Here $\rho$ is the usual frame residue modulo three.
The nine finite affine translations of the descended cubic
frame allowed by its exact frame equation give all three
values of $u$, for each of the two values of $b=\Psi_3/2\bmod2$.
This is a necessary finite list, not a lifting assertion for
each entry.

For every value in (D23.F27.1), the periodic cube equation
of D23.C17 holds with either possible residual sign
\[
          \eta_3(D_h)=(0,(-I)^e),\qquad e\in\mathbb F_2.
\tag{D23.F27.2}
\]
Thus this non-Ferrari periodicity test does not exclude the
sign $e=1$.  The affine periodic coinvariant remains nonzero
for both signs.  No smaller comparison gate or vanishing of
$e_3$ is inferred from this calculation.

\begin{proof}
We first identify the marked exceptional pair, not just its
abstract finite group.  In the cubic detector of D23.P3, let
$\eta$ be the affine representation over $\mathbb F_9$ with
\[
 \eta(s)=(0,iS),\qquad \eta(r)=(e_1,R),\quad
 S=\begin{pmatrix}0&1\\-1&0\end{pmatrix},\quad R=-UL.
\]
On the index-two subgroup $E=\langle r,srs\rangle$, all
linear parts have entries in $\mathbb F_3$, so real-part
projection of translations is a homomorphism.  Follow it
by the origin change $t=(0,2)^{\mathsf t}$, as in D23.P10
and D23.C14.  The resulting three marked elements are
\[
 x=\left(0,\begin{pmatrix}2&1\\0&2\end{pmatrix}\right),\quad
 r=\left(\binom22,\begin{pmatrix}0&2\\1&2\end{pmatrix}\right),\quad
 v=srs=\left(\binom11,\begin{pmatrix}2&2\\1&0\end{pmatrix}\right).
\]
Conjugation by the linear matrix
\[
                         H=\begin{pmatrix}2&0\\1&2\end{pmatrix}
\tag{D23.F27.3}
\]
sends the pair $(x,v)$ to
\[
 \eta_3(\delta)=\left(0,\begin{pmatrix}0&1\\2&1\end{pmatrix}\right),
 \qquad
 \eta_3(K)=\left(-e_1,\begin{pmatrix}1&2\\0&1\end{pmatrix}\right).
\tag{D23.F27.4}
\]
These are literal equalities, checked on both inputs.
In the original modular group $\tau=sr$ fixes $x=\tau^{-2}$
by conjugation and sends $r$ to $v$.  Consequently
\[
                    f(x,v)=\tau C_f\tau^{-1},\qquad C_f=f(x,r).
\tag{D23.F27.5}
\]
The conjugation in this equation must be evaluated before real
projection: $\tau$ is outside $E$, and that projection is not
being extended to the whole modular group.

Here are the full finite frame constraints and their evaluation.
Let $A_E$ be D23.F22's descended pure cubic frame and put
$C_\rho=U^2M_\rho U^{-2}$, with
$M_\rho=\left(\begin{smallmatrix}1&\rho\\\rho&1+\rho^2\end{smallmatrix}\right)$.
D23.F22 gives its linear part
\[
                  A_0=(-I)^b C_\rho^{-1}.
\]
D23.P3 gives the complete affine value of the descended pair
coordinate $f_E$, using $\rho(D_Eh)=\rho(h)$ and
$a_3(D_Eh)=b$.  Thus the translation $t_E$ of $\eta(A_E)$
satisfies the exact affine equation
\[
       \eta(A_E)\eta(r)\eta(A_E)^{-1}
                         =\eta(f_E)^{-1}\eta(r).
\tag{D23.F27.6}
\]
Solving its two linear translation equations over $\mathbb F_9$
gives the following complete parametrizations, with $z\in\mathbb F_9$.
The last column is the first translation coordinate after
forming $C_f=A_E^{-1}x^{-3b}$, applying (D23.F27.5), real
projection with the stated origin shift, and conjugation by $H$.
Write $\operatorname{Im}(a+bi)=b\in\mathbb F_3$.
\[
\begin{array}{cc|c|c}
 \rho&b&t_E&u\\\hline
 0&0&(z,-z)&\operatorname{Im}z\\
 0&1&(z,1+i-z)&2-\operatorname{Im}z\\
 1&0&(z,2+2i+z)&-\operatorname{Im}z\\
 1&1&(z,2+i+z)&1+\operatorname{Im}z\\
 2&0&(2+2i,z)&1-\operatorname{Im}z\\
 2&1&(1,z)&2+\operatorname{Im}z.
\end{array}
\tag{D23.F27.7}
\]
In every row the second resulting translation coordinate is
zero and the linear part is $N_\rho=C_\rho^{-1}$.
For an explicit check of the powers appearing here, D23.P3's
affine coordinates are
\[
 \eta(\tau)=((0,-i),iU^{-1}),\qquad
 \eta(x)=((1+i,1+2i),-U^2),\qquad
 \eta(x)^3=((-i,1+2i),-I).
\]
Substituting these coordinates in the affine multiplication
law verifies every row of (D23.F27.7).  In particular no central
linear factor has been discarded as an affine translation.
The table proves (D23.F27.1) and explains why its set is
independent of $b$.

Since $D_h$ is evaluated in the linear first-three-strand
subgroup, D23.C17.2 gives exactly (D23.F27.2).  Put
$\epsilon=(-1)^e$.  The translation $u e_1$ in (D23.F27.1)
commutes with $\eta_3(K)$ because $U^2e_1=e_1$.  Direct affine
multiplication therefore gives
\[
 \eta_3(\delta D_hg^{-1}KgD_h^{-1})
             =(\epsilon v_\rho,L_\rho),
\]
where
\[
\begin{array}{c|c|c|c}
 \rho&L_\rho&v_\rho&\ell_\rho\\\hline
 0&\left(\begin{smallmatrix}0&1\\2&2\end{smallmatrix}\right)
       &(0,1)^{\mathsf t}&(2,1)\\
 1&\left(\begin{smallmatrix}1&1\\0&1\end{smallmatrix}\right)
       &(2,2)^{\mathsf t}&(0,2)\\
 2&\left(\begin{smallmatrix}1&0\\2&1\end{smallmatrix}\right)
       &(1,0)^{\mathsf t}&(1,0).
\end{array}
\tag{D23.F27.8}
\]
Each $L_\rho$ is a nonidentity unipotent with
$(L_\rho-I)^2=0$.  In characteristic three this implies
$I+L_\rho+L_\rho^2=0$.  Hence every affine lift has cube one,
for either sign of its translation.  The last column verifies
$\ell_\rho L_\rho=\ell_\rho$ and
$\ell_\rho v_\rho=1$.  Its affine periodic coinvariant is
therefore $\epsilon$, nonzero for either sign.
Conjugation by $(0,-I)$ interchanges these two nonzero classes;
requiring ordinary affine conjugacy does not distinguish them.

There is also a simpler general reason the displayed residual
sign cannot enter this test directly.  In this quotient
$\eta_3(D_h)$ commutes with $\eta_3(\delta)$, so the whole
periodic product is conjugate by $\eta_3(D_h)$ to
$\eta_3(\delta g^{-1}Kg)$.  The explicit calculation above
checks that the possible actual exceptional values do satisfy
the cube equation, rather than merely observing that the
residual sign drops out of it.

The certificate \path{agents/five_level_six/orbital_affine216.py}
rechecks the two exact marked conjugacies, retains the
$\mathbb F_9$ conjugation by $\tau$, enumerates all $81$
translation candidates in each of the six frame equations,
and checks all $108$ remaining signed frame cases.  Its
record is \path{orbital_affine216.json}.  The full joint
orbital identities contain additional correlated evaluations;
this proposition only establishes the stated obstruction for
the periodic residual in this specific quotient.
\end{proof}
